\section*{第 \ref{ch:b1:nernst} 章——氧化还原平衡与能斯特方程}

\begin{solution}{exo:b1:nernst:1}
N 在 \ce{NH3} 中为 $-$III；在 \ce{NO3-} 中为 $+$V。S 在 \ce{SO4^2-} 中为 $+$VI。O 在
\ce{H2O2} 中为 $-$I。Cr 在 \ce{Cr2O7^2-} 中为 $+$VI。Cl 在 \ce{ClO-} 中为 $+$I。Fe 在
\ce{Fe3O4} 中 Fe 的平均\omterm{def:b1:nernst:oxidation-number}{氧化数}为 $+8/3$（一个铁(II)，两个铁(III)）。C 在
\ce{CH3OH} 中，$x + 4(+1) + (-2) = 0$，即 $-$II。
\end{solution}

\begin{solution}{exo:b1:nernst:2}
\ce{Cr2O7^2- + 14H+ + 6e- -> 2Cr^3+ + 7H2O};
\ce{H2O2 + 2H+ + 2e- -> 2H2O};
\ce{O2 + 2H+ + 2e- -> H2O2};
\ce{SO4^2- + 4H+ + 2e- -> SO2 + 2H2O}.
\end{solution}

\begin{solution}{exo:b1:nernst:3}
$E = -0.76 + 0.030\log[\ce{Zn^2+}]$; $E = 1.36 + 0.030\log(p(\ce{Cl2})/[\ce{Cl-}]^2)$;
$E = 1.23 + 0.015\log(p(\ce{O2})[\ce{H+}]^4)$; $E = 0.27 - 0.059\log[\ce{Cl-}]$
（压强以 bar 为单位）。在 \qty{0.010}{mol/L} 硫酸锌溶液中：
$E = -0.76 + 0.030 \times (-2) = -0.82$~V.
\end{solution}

\begin{solution}{exo:b1:nernst:4}
$E(\ce{Zn}) = -0.76 + 0.030\log 0.10 = -0.79$~V, $E(\ce{Cu}) = 0.34 +
0.030\log 0.010 = 0.28$~V。电压为 $0.28 - (-0.79) = 1.07$~V。锌
电势较低，为负极，是发生氧化的
\omterm{def:b1:nernst:cell}{阳极}。
\end{solution}

\begin{solution}{exo:b1:nernst:5}
\ce{MnO4- + 5Fe^2+ + 8H+ -> Mn^2+ + 5Fe^3+ + 4H2O}, $n = 5$,
$\log K = 5(1.51 - 0.77)/0.059 = 63$：定量且快速，是经典的
滴定（高锰酸根自身即为指示剂）。
\end{solution}

\begin{solution}{exo:b1:nernst:6}
\omterm{def:b1:nernst:electrode-potential}{标准电极电势} $E^\circ$ 低于 0 的金属可以反应：Mg、Zn、Fe、Ni、Pb（铅只
缓慢反应，电势差较小）；Cu 和 Ag 则不反应。\ce{Mg + 2H+ -> Mg^2+ + H2}, $\log K = 2 \times 2.36/0.059 = 80$；
\ce{Fe + 2H+ -> Fe^2+ + H2}, $\log K = 2 \times 0.41/0.059 = 13.9$.
\end{solution}

\begin{solution}{exo:b1:nernst:7}
$n = 2$，$\log K = 2(0.53 - 0.02)/0.059 = 17$：\omterm{def:b1:extent-q-and-k:quantitative}{定量反应}。
\end{solution}

\begin{solution}{exo:b1:nernst:8}
$3E^\circ(\ce{Fe^3+}/\ce{Fe}) = 2(-0.41) + 0.77$, $E^\circ = -0.02$~V.
\ce{2Fe^3+ + Fe -> 3Fe^2+}：\omterm{def:b1:nernst:couple}{氧化剂} \ce{Fe^3+}（0.77）高于\omterm{def:b1:nernst:couple}{还原剂}
\ce{Fe}（$-0.41$），$n = 2$，$\log K = 2(0.77 + 0.41)/0.059 = 40$。铁
屑可防止铁(II)溶液变为铁(III)。
\end{solution}

\begin{solution}{exo:b1:nernst:9}
\ce{O2 + 4H+ + 4e- -> 2H2O}: $E = 1.23 + \frac{0.059}{4}\log [\ce{H+}]^4 =
1.23 - 0.059\,\mathrm{pH}$. \ce{2H+ + 2e- -> H2}: $E = \frac{0.059}{2}\log
[\ce{H+}]^2 = -0.059\,\mathrm{pH}$。两者之差在任何 pH 下均为 $1.23$~V：
燃料电池的电压与 pH 无关。
\end{solution}

\begin{solution}{exo:b1:nernst:10}
$E^\circ(\ce{Cu+}/\ce{Cu}) = 0.52 > E^\circ(\ce{Cu^2+}/\ce{Cu+}) = 0.16$：上电对的
\omterm{def:b1:nernst:couple}{氧化剂} \ce{Cu+} 与
下电对的\omterm{def:b1:nernst:couple}{还原剂} \ce{Cu+} 发生反应。$\log K = (0.52 - 0.16)/0.059 = 6.1$，$K =
[\ce{Cu^2+}]/[\ce{Cu+}]^2$，因此 $[\ce{Cu+}] = \sqrt{0.10/10^{6.1}} =
\qty{2.8e-4}{mol/L}$.
\end{solution}

\begin{solution}{exo:b1:nernst:11}
$U = 0.059\log(0.10/\num{1.0e-3}) = 0.118$~V。正极为
浓溶液中的银（此处 \ce{Ag+} 被还原）；在
稀溶液中的银被氧化。浓溶液变得
更稀，稀溶液则变得更浓；电压下降，
并在两浓度相等时降为零。
\end{solution}

\begin{solution}{exo:b1:nernst:12}
\ce{H2O2} 是 \ce{H2O2}/\ce{H2O}（1.76~V）的\omterm{def:b1:nernst:couple}{氧化剂}，又
是 \ce{O2}/\ce{H2O2}（0.69~V）的\omterm{def:b1:nernst:couple}{还原剂}：它与自身反应，\ce{2H2O2 -> 2H2O + O2}，
$\log K = 2(1.76 - 0.69)/0.059 = 36$。该反应热力学上有利，但
没有\omterm{def:b1:elementary-steps:catalyst}{催化剂}时非常缓慢（见\cref{ch:b1:elementary-steps}）；瓶子洁净，
低温、避光保存，可使其保持缓慢。
\end{solution}

\begin{solution}{pb:b1:nernst:1}
\textbf{1.} Ag：0，$+$I，$+$I；Cl 在两者中均为 $-$I。
\textbf{2.} \ce{AgCl(s) + e- -> Ag(s) + Cl-}.
\textbf{3.} $E = E^\circ(\ce{AgCl}/\ce{Ag}) - 0.059\log[\ce{Cl-}]$.
\textbf{4.} 两种固体的\omterm{def:b1:extent-q-and-k:activity}{活度}均为 1；只有氯离子
存在于溶液中。
\textbf{5.} $E^\circ = 0.22$~V.
\textbf{6.} $0.22 + 0.118 = 0.34$~V.
\textbf{7.} \ce{Pt} $|$ \ce{H2}（\qty{1}{bar}）$|$ \ce{H+}（\omterm{def:b1:extent-q-and-k:activity}{活度}为 1）
$\|$ \ce{Cl-} (\qty{0.10}{mol/L}) $|$ \ce{AgCl} $|$ \ce{Ag}.
\textbf{8.} 银电极：$0.22 + 0.059 = 0.28$~V $> 0$。
\textbf{9.} $0.28$~V.
\textbf{10.} \ce{2AgCl + H2 -> 2Ag + 2H+ + 2Cl-}.
\textbf{11.} $\log K = 2 \times 0.22/0.059 = 7.5$.
\textbf{12.} 由铂（负极，\ce{H2} 在此被氧化）
流向银。
\textbf{13.} $0.27 - 0.059\log 1.0 = 0.27$~V.
\textbf{14.} $E(\ce{Ag}) = 0.22 + 0.177 = 0.40$~V $> 0.27$：银
电极为正极，$U = 0.13$~V。
\textbf{15.} $U = 0.22 - 0.059\log[\ce{Cl-}] - 0.27 = -0.05 -
0.059\log[\ce{Cl-}]$（单位为伏特）。
\textbf{16.} $\log[\ce{Cl-}] = -(0.115 + 0.05)/0.059 = -2.80$,
$[\ce{Cl-}] = \qty{1.6e-3}{mol/L}$，即 \qty{57}{mg/L}。
\textbf{17.} $\Delta c/c = \ln 10 \times 0.001/0.059 = 0.039$：约
\qty{4}{\percent}，每毫伏。
\textbf{18.} 当样品中的氯离子与
氯化银自身所释放者相当时，$\sqrt{K_s} = \qty{1.3e-5}{mol/L}$；
电极宜用于远高于此值的场合，约自 \qty{1e-4}{mol/L} 起。
\textbf{19.} $E = 0.80 + 0.059\log[\ce{Ag+}]$.
\textbf{20.} $[\ce{Ag+}] = K_s/[\ce{Cl-}]$.
\textbf{21.} $E = 0.80 + 0.059\log K_s - 0.059\log[\ce{Cl-}] = 0.80 -
0.059\,\mathrm{p}K_s - 0.059\log[\ce{Cl-}]$.
\textbf{22.} $E^\circ(\ce{AgCl}/\ce{Ag}) = E^\circ(\ce{Ag+}/\ce{Ag}) -
0.059\,\mathrm{p}K_s$ (\cref{prop:b1:nernst:second-kind}).
\textbf{23.} $0.80 - 0.059 \times 9.75 = 0.225$~V，与
表中数值（0.22~V）相符，由吉布斯自由能独立算出。
\textbf{24.} $E^\circ(\ce{AgCl}/\ce{Ag}) = $ \textbf{0.22~V}.
\end{solution}
